\subsection{UI Overlay}

Rather than hard-coding a fixed HUD, the overlay is exposed as a
general-purpose compositing layer. This was a deliberate design choice to
position the system as a reusable game engine rather than a
single-game raycaster. Any application layer can define its own
screen regions as the hardware enforces no assumptions about what is
drawn where.

\vspace{0.5em}

The overlay pipeline inserts itself in the AXI-Stream between the
raycaster output and the HDMI encoder. At each pixel clock cycle it
receives the current beam coordinates $(x, y)$ from the video timing
controller and selects the output colour according to a strict
priority hierarchy:

\begin{keyequation}
\[
\text{pixel}_{\text{out}} =
\begin{cases}
C_{\text{sprite}}(x,y) & \text{if } (x,y) \in R_{\text{sprite}} \\
C_{\text{glyph}}(x,y)  & \text{if } (x,y) \in R_{\text{text}}   \\
C_{\text{bar}}         & \text{if } (x,y) \in R_{\text{bar}}    \\
C_{\text{ray}}(x,y)    & \text{otherwise}
\end{cases}
\]
\end{keyequation}

Each region $R$ is defined by a pair of axis-aligned bounding boxes
written to AXI-Lite registers by the PS at startup. Because the
decision resolves in a single clock cycle, with no memory transactions
on the critical path, the overlay adds \textbf{zero pipeline stages}
and consumes no additional BRAM beyond the glyph ROM.

\vspace{0.5em}

Text rendering uses a $760 \times 8$ BRAM-backed glyph ROM covering
the full 95-character printable ASCII set. For a character at grid
position $(c_x, c_y)$ the ROM address is:

\begin{keyequation}
\[
\text{addr} = \bigl(\text{ASCII}(c) - 32\bigr) \times 8
              + \bigl(y - c_y \cdot 8\bigr)
\]
\end{keyequation}

and the pixel colour is gated by bit $\bigl(x - c_x \cdot 8\bigr)$
of the returned byte, giving a 1-cycle-per-pixel throughput with no
stalls.

\vspace{0.5em}

A proof-of-concept sprite layer was implemented alongside the text
renderer to validate the extension path: a $16 \times 16$ icon is
loaded into a dedicated BRAM region at boot and composited at a
PS-configurable screen coordinate. The same mechanism generalises to
arbitrary sprite sheets where a wider AXI-Lite register map could expose
multiple independently-positioned sprites at runtime without modifying
the core pipeline.

\begin{figure}[H]
    \centering
    % TODO: screenshot of game running with HUD overlaid, or pipeline
    % block diagram showing pixel-mux between raycaster / glyph ROM /
    % bar logic / sprite BRAM
    \fbox{\parbox{0.85\columnwidth}{\centering\vspace{2em}
        \textit{[Figure placeholder: annotated screenshot showing an example UI]}
    \vspace{2em}}}
    \caption{UI overlay: pixel-clock-domain compositing over raycasted frame.}
    \label{fig:ui-overlay}
\end{figure}

\paragraph{Resource cost:}
The entire overlay (region comparators, glyph ROM, sprite BRAM, and
pixel mux) consumes ~180 LUTs and 1 BRAM tile
(36\,Kb). This is substantially cheaper than a software alternative:
blitting a $320 \times 240$ framebuffer from the PS would transfer
$\sim$75\,kB per frame, consuming $\sim$36\,MB/s of AXI bandwidth at
60\,fps and competing directly with the network and game-loop threads.
The PL overlay eliminates that cost entirely.